% Source PDF: source/普通高考/2007/2007广东文.pdf, page 1, question 7, figure 2.
% The decision diamond is intentionally blank in the source; it is the answer blank.
% Topology: 开始 -> 输入 A_1,...,A_{10} -> s=0,i=4 -> blank decision.
% 是 -> s=s+A_i -> i=i+1 -> shared entry above the decision; 否 -> 输出 s -> 结束.

\begin{tikzpicture}[
  x=1cm,y=1cm,baseline,>=Stealth,
  line width=0.45pt,line cap=butt,line join=miter,
  every node/.style={anchor=center,align=center,font=\normalsize,
    execute at begin node={\setlength{\parindent}{0pt}\noindent}},
  flowbox/.style={draw,anchor=center,align=center,inner xsep=4pt,inner ysep=2pt,
    execute at begin node={\setlength{\parindent}{0pt}\noindent}},
  terminal/.style={draw,anchor=center,align=center,inner xsep=4pt,inner ysep=2pt,
    execute at begin node={\setlength{\parindent}{0pt}\noindent},
    rounded corners=0.16cm,minimum width=1.35cm,minimum height=0.68cm},
  process/.style={draw,anchor=center,align=center,inner xsep=4pt,inner ysep=2pt,
    execute at begin node={\setlength{\parindent}{0pt}\noindent},
    minimum height=0.68cm},
  decision/.style={draw,anchor=center,align=center,inner xsep=4pt,inner ysep=2pt,
    execute at begin node={\setlength{\parindent}{0pt}\noindent},
    diamond,aspect=2.85,minimum width=3.20cm,minimum height=1.05cm},
  io/.style={draw,anchor=center,align=center,inner xsep=4pt,inner ysep=2pt,
    execute at begin node={\setlength{\parindent}{0pt}\noindent},
    trapezium,trapezium left angle=72,trapezium right angle=108,
    minimum height=0.76cm}
]
  \node[terminal] (start) at (0,9.20) {开始};
  \node[io,minimum width=4.85cm] (input) at (0,7.75) {输入 $A_1,A_2,\ldots,A_{10}$};
  \node[process,minimum width=2.65cm] (init) at (0,6.25) {$s=0,i=4$};
  \node[decision] (test) at (0,3.90) {};
  \node[process,minimum width=2.55cm] (sum) at (3.30,3.90) {$s=s+A_i$};
  \node[process,minimum width=2.25cm] (inc) at (3.30,5.65) {$i=i+1$};
  \node[io,minimum width=1.90cm] (output) at (0,2.10) {输出 $s$};
  \node[terminal] (finish) at (0,0.60) {结束};

  \draw[->] (start.south) -- (input.north);
  \draw[->] (input.south) -- (init.north);
  \coordinate (merge) at (0,4.90);
  \coordinate (loopjoin) at (0,5.65);
  \draw (init.south) -- (merge);
  \draw[->] (merge) -- (test.north);
  \draw[->] (test.east) -- node[above,inner sep=2pt] {是} (sum.west);
  \draw[->] (sum.north) -- (inc.south);
  % Return path merges into the main stem; the arrowhead marks the
  % leftward flow into the shared stem before the decision.
  \draw[->] (inc.west) -- (loopjoin);
  \draw[->] (test.south) -- node[right,inner sep=2pt] {否} (output.north);
  \draw[->] (output.south) -- (finish.north);
\end{tikzpicture}
